% Hikaru's A adaptation of Desmond's B placement ideas.
\subsection{\expt{a111}: Colocation on join keys}
Experiment a111 retained a102's indexes and adapted the B team's placement:
\begin{itemize}
\item Distribute \code{inventory} by \code{product_id} and colocate it with
\code{products}.
\item Distribute \code{order_items} by \code{order_id} and colocate it with
\code{orders}.
\item Add the corresponding product identifier to T7's inventory-update
routing predicates.
\end{itemize}
The implementation retains the selected inventory identifier and does not
assume that \code{inventory_id} equals \code{product_id}.

Matching rows are available together for the relevant joins. This does not
mean every whole query accesses one worker: an order-history query may involve
multiple shards, while matching orders and items can be joined within colocated
shard pairs \citep{citus-distribution}. The scale-1, four-client diagnostic,
job 977617, achieved 89.917 TPS with successful phases, no measured failures,
and passing validation.
